\section{Definitions}\label{sec:tsps-definitions}

Structure preserving signatures sign messages that lie in pairing friendly
groups, rather than arbitrary strings, and the corresponding verification
algorithm checks equations comprising products of pairings. A threshold
structure preserving signature scheme distributes the signing key of such a
scheme across $n$ signers, any $t+1$ of whom can jointly produce a signature,
and no $t$ of whom learn anything about the key.

The signers are the only participants in the protocol. They communicate over
authenticated point to point channels, together with a broadcast channel, and
there is no trusted dealer: the signing key is never held in full by any party,
at any point, including during key generation. A verifier holds only the public
key, and needs to know neither $n$, nor $t$, nor which signers took part.
Signers may be arbitrarily malicious, deviating from the protocol at will, and
may collude to forge signatures or to recover the signing key. We consider
adaptive corruption: the adversary chooses which signers to corrupt during
protocol execution, subject to the bound $|\Cor| \leq t$, and we write $\Cor$
for the set of corrupted signers and $\mathcal{H}$ for its complement.

We write group operations multiplicatively throughout, with $\gen$ and $\ggen$
the fixed generators of $\group{1}$ and $\group{2}$ respectively.

\subsection{Structure preserving signatures}

A structure preserving signature scheme is defined by the following algorithms:

\begin{description}

  \item[$(\pk, \sk) \sample \SPSgen(\secparam, \ell)$ :] The key generation
algorithm, on input security parameter $\secparam$ and length $\ell$, outputs
key pair $(\pk, \sk)$. Public keys are group elements.

  \item[$\SPSsig \sample \SPSsign(\sk, \vec{\msg})$ :] The signing algorithm, on
input secret key $\sk$ and message $\vec{\msg}$, outputs signature $\SPSsig$.
Messages and signatures are vectors of group elements, messages from the space
$\group{1}^\ell$, and signatures in $\group{1}^{\mathrm{poly}(\secparam)} \times
\group{2}^{\mathrm{poly}(\secparam)}$.

  \item[$\{0,1\} \leftarrow \SPSverify(\pk, \vec{\msg}, \SPSsig)$ :] The
verification algorithm, on input public key $\pk$, message $\vec{\msg}$, and
signature $\SPSsig$, returns 0 or 1. Verification consists only of evaluating
pairing product equations.

\end{description}

A structure preserving signature scheme satisfies correctness if, for all
security parameters $\secparam$, all $(\pk, \sk) \sample \SPSgen(\secparam,
\ell)$, and all $\vec{\msg} \in \group{1}^\ell$, the following holds:
%
\begin{equation*}
%
  \SPSverify(\pk, \vec{\msg}, \SPSsign(\sk, \vec{\msg})) = 1.
%
\end{equation*}
%
As in~\Cref{sec:prelims-encryption}, we write $\SPSsign(\sk, \vec{\msg}; \rho)$
where a security argument needs to name the signing randomness explicitly.
%

\paragraph{Unforgeability.} For a PPT adversary $\adv$, define advantage as
follows:
%
\begin{equation*}
%
  \Adv_{\SPS, \adv}^{\text{UF-CMA}}(\secparam) = \Pr\left[ \mathbf{G}_{\SPS,
\adv}^{\text{UF-CMA}}(\secparam) = 1\right].
%
\end{equation*}
%
Here, game $\mathbf{G}_{\SPS,\adv}^{\text{UF-CMA}}$ and signing oracle
$\TSPSsigno$ are given in Experiment~\ref{exp:sps-uf-cma}. A structure
preserving signature scheme achieves UF-CMA security if, for every PPT adversary
$\adv$, the following holds:
%
\begin{equation*}
%
  \Adv_{\SPS, \adv}^{\text{UF-CMA}}(\secparam) \leq \negl(\secparam).
%
\end{equation*}
%
Experiment~\ref{exp:sps-uf-cma} excludes $\vec{\msg}^*$ from the set $S$ of
queried messages, so that a forgery must be on a message the adversary never
asked about. This is existential unforgeability, EUF-CMA. Excluding instead the
pair $(\vec{\msg}^*, \SPSsig^*)$, so that a second, different signature on an
already queried message also counts as a forgery, gives strong unforgeability,
sUF-CMA, with advantage $\Adv_{\SPS, \adv}^{\text{sUF-CMA}}(\secparam)$. The
sUF-CMA winning condition is the more permissive of the two, so sUF-CMA is the
stronger notion and implies EUF-CMA.

EUF-CMA is the notion we assume of the underlying scheme throughout this
chapter, and we take the weaker of the two deliberately rather than for
convenience. Our construction is generic: it takes an arbitrary structure
preserving signature scheme and thresholdises it, so the weaker the assumption
we place on that scheme, the wider the class of schemes the construction covers.
Assuming only EUF-CMA costs us nothing, because the threshold notion we target
screens forgeries by message rather than by message and signature together. An
adversary against the threshold scheme wins only on a $\vec{\msg}^*$ for which
too few honest signers issued a partial signature.  The reduction of
\Cref{sec:tsps-proof-1} is arranged so that it queries its own oracle only on
messages falling outside that condition. The forgery it forwards is therefore
always on a message it never queried, which is exactly what EUF-CMA asks for,
and at no point does it need to distinguish two signatures on the same message.
Strong unforgeability is consequently never used, and requiring it would only
narrow the class of schemes the construction applies to.

Nothing is lost at the threshold level either. Should an application require the
strong guarantee, the same reduction gives the strong threshold analogue
whenever the underlying scheme is sUF-CMA. The only change is to the winning
condition carried across, and by \Cref{lem:tsps-transparency} a combined
signature is distributed exactly as an ordinary one. The threshold scheme
inherits whichever of the two notions the underlying scheme achieves, and we
state our results under the weaker hypothesis so that they apply to both.

\begin{experimentbox}{Unforgeability under chosen message attack for structure
preserving signatures}{sps-uf-cma}

  \begin{enumerate}

    \item The challenger generates $(\pk, \sk) \gets \SPSgen(\secparam, \ell)$
and sends $\pk$ to $\adv$.

    \item The adversary is given oracle access to $\TSPSsigno$, which on input
$\vec{\msg}$, computes $\SPSsig \gets \SPSsign(\sk, \vec{\msg})$, and if
$\SPSsig \neq \perp$, adds $\vec{\msg}$ to $S$, before returning $\SPSsig$.

    \item The adversary outputs $(\vec{\msg}^*, \SPSsig^*)$.

  \end{enumerate}

  $\adv$ wins the experiment if $\SPSverify(\pk, \vec{\msg}^*, \SPSsig^*) = 1$
and $\vec{\msg}^* \notin S$.

\end{experimentbox}

\subsection{Threshold structure preserving signatures}

A threshold structure preserving signature scheme is a structure preserving
signature scheme whose signing key is distributed across multiple parties.
Because our constructions perform all of their interaction ahead of the message
being known, we separate signing into an offline and an online algorithm. We
call a scheme partially non-interactive when $\algo{OnSign}$ requires no
communication between signers while $\algo{OffSign}$ may. The scheme is defined
by algorithms $\TSPSgen$, $\algo{OffSign}$, $\algo{OnSign}$, $\TSPScombine$, and
$\TSPSverify$ as follows.

\begin{description}

  \item[{$(\pk, \shr{\sk}) \sample \TSPSgen(\secparam, \ell, t, n)$ :}] The key
generation algorithm takes a triple of integers $(\ell, t, n > t)$, where $\ell$
is the length of the vectors to be signed, $t$ is an upper bound on the number
of corrupted parties, and $n$ is the total number of threshold signers. It
returns a public key $\pk$ and $n$ shares $\shr{\sk}$ such that any $t+1$ shares
reconstruct the secret key underlying $\pk$, and any $t$ shares are independent
of it. $\pk$ consists of group elements, and $\sk$ of elements of $\Zp$ and/or
group elements.

  \item[{$\shr{\state} \sample \algo{OffSign}(\shr{\sk})$ :}] The offline
signing algorithm is run jointly by the signers, before the message is known. It
takes the secret key shares and outputs a share of pre-processed state
$\shr{\state}$, comprising those components of the signature that do not depend
on the message together with the material $\algo{OnSign}$ consumes.  Each
invocation produces state for exactly one future signature, and that state is
single-use.

  \item[{$\shr{\SPSsig} \leftarrow \algo{OnSign}(\shr{\sk}, \shr{\state},
\vec{\msg})$ :}] The online signing algorithm is run locally by each signer. It
takes that signer's key share, its share of a previously unused pre-processed
state, and a message $\vec{\msg} \in \group{1}^\ell$, and outputs a signature
share $\shr{\SPSsig}$ with elements in $\group{1}$ and $\group{2}$.

  \item[{$\SPSsig \leftarrow \TSPScombine(\shr{\SPSsig})$ :}] The combining
algorithm is non-interactive, and takes $t+1$ signature shares $\shr{\SPSsig}$
to output signature $\SPSsig$.

  \item[$\{0,1\} \leftarrow \TSPSverify(\pk, \vec{\msg}, \SPSsig)$ :] The
verification algorithm takes public key $\pk$, message $\vec{\msg}$, and
signature $\SPSsig$, and returns 0 or 1. $\TSPSverify$ is equal to $\SPSverify$.

\end{description}

\paragraph{Correctness.} A threshold structure preserving signature scheme
satisfies correctness if, for all $\secparam$, all $(\ell, t, n > t)$, all
$(\pk, \shr{\sk}) \sample \TSPSgen(\secparam, \ell, t, n)$, all $\shr{\state}
\sample \algo{OffSign}(\shr{\sk})$, all $\vec{\msg} \in \group{1}^\ell$, and
every subset $Q \subseteq [1,n]$ with $|Q| = t+1$, the signature obtained by
running $\algo{OnSign}$ at each signer in $Q$ and combining the results
verifies:
%
\begin{equation*}
%
  \SPSverify\left(\pk, \vec{\msg}, \TSPScombine\left(\left\{
\algo{OnSign}(\sk_i, \state_i, \vec{\msg}) \right\}_{i \in Q}\right)\right) = 1.
%
\end{equation*}
%
Correctness is required to hold for every such $Q$, so that the signature does
not depend on which $t+1$ signers took part.

\paragraph{Unforgeability.} Threshold signatures admit a hierarchy of
unforgeability notions~\cite{C:BCKMTZ22}. We rely on adaptive TS-UF-1 security,
in which the adversary has an offline oracle $\mathcal{O}_{\algo{Off}}$, a
partial signature oracle $\TSPSsigno$, and a corruption oracle $\corrupto$, and
wins with $(\vec{\msg}^*, \SPSsig^*)$ provided it obtained partial signatures on
$\vec{\msg}^*$ from fewer than $t + 1 - |\Cor|$ honest signers. That is, the
partial signatures it holds on $\vec{\msg}^*$, together with the key shares of
the signers it corrupted, must fall short of the $t+1$ that reconstruct a
signature. The stricter TS-UF-0 disqualifies the adversary as soon as
$\TSPSsigno$ is queried on $\vec{\msg}^*$ even once. That is too strong for our
purposes. our online phase is local, so a signer that has issued a single
partial signature has revealed nothing further.

The offline oracle is what makes the notion match our syntax. All interaction
between signers happens in $\algo{OffSign}$, so an adversary that could not
invoke it would never see the protocol messages in which that interaction is
carried out. $\mathcal{O}_{\algo{Off}}$ therefore runs the offline phase with
$\adv$ playing the corrupted signers, and returns the resulting view. Each
session identifier may be used once by $\mathcal{O}_{\algo{Off}}$ and once by
$\TSPSsigno$, reflecting the single-use requirement on pre-processed state:
reusing a state across two messages would reuse the signing randomness, which
neither underlying signature scheme tolerates.

Let $\mathrm{TSPS} = (\TSPSgen, \algo{OffSign}, \algo{OnSign}, \TSPScombine,
\SPSverify)$ be an $(n, t)$-TSPS scheme over message space $\mathcal{M}$. The
advantage of a PPT adversary $\adv$ is defined as follows:
%
\begin{equation*}
%
  \Adv_{\TSPS, \adv}^{\text{adp-TS-UF-1}}(\secparam) = \Pr\left[
\mathbf{G}_{\TSPS, \adv}^{\text{adp-TS-UF-1}}(\secparam) = 1 \right].
%
\end{equation*}
%
Here, game $\mathbf{G}_{\TSPS, \adv}^{\text{adp-TS-UF-1}}$ and its oracles
$\mathcal{O}_{\algo{Off}}$, $\TSPSsigno$ and $\corrupto$ are given in
Experiment~\ref{exp:tsps-adp-ts-uf-1}. A threshold structure preserving
signature scheme achieves adaptive TS-UF-1 security if, for every PPT adversary
$\adv$, the following holds:
%
\begin{equation*}
%
  \Adv_{\TSPS, \adv}^{\text{adp-TS-UF-1}}(\secparam) \leq \negl(\secparam).
%
\end{equation*}

\begin{experimentbox}{Adaptive unforgeability under chosen message attack for
threshold structure preserving signatures (adp-TS-UF-1)}{tsps-adp-ts-uf-1}

  \begin{enumerate}

    \item The adversary outputs $(n, t, \Cor)$, and the challenger sets
$\mathcal{H} := [1,n] \backslash \Cor$.

    \item The challenger generates $(\pk, \shr{\sk}) \gets \TSPSgen(\secparam,
\ell, t, n)$, and sends $\pk$ and $\{\sk_i\}_{i \in \Cor}$ to $\adv$.

    \item The adversary is given oracle access to $\mathcal{O}_{\algo{Off}}$,
which on input a fresh session identifier $\mathsf{sid}$, runs $\algo{OffSign}$
among all $n$ signers with $\adv$ acting on behalf of $\Cor$, stores the
resulting shares $\shr{\state}_{\mathsf{sid}}$, and returns the view of the
corrupted signers. It returns $\perp$ if $\mathsf{sid}$ has been used before.

    \item The adversary is given oracle access to $\TSPSsigno$, which on input
$(i, \mathsf{sid}, \vec{\msg})$, asserts $\vec{\msg} \in \mathcal{M}$, $i \in
\mathcal{H}$, and that $\mathsf{sid}$ was initialised by
$\mathcal{O}_{\algo{Off}}$ and not yet used for signing. It then computes
$\SPSsig_i \gets \algo{OnSign}(\sk_i, \state_{\mathsf{sid},i}, \vec{\msg})$, and
if $\SPSsig_i \neq \perp$, adds $i$ to $S(\vec{\msg})$ before returning
$\SPSsig_i$.

    \item The adversary is given oracle access to $\corrupto$, which on input
$k$, returns $\perp$ if $k \in \Cor$ or $|\Cor| = t$, and otherwise adds $k$ to
$\Cor$, removes $k$ from $\mathcal{H}$, and returns $\sk_k$ together with every
$\state_{\mathsf{sid},k}$ held by signer $k$.

    \item The adversary outputs $(\vec{\msg}^*, \SPSsig^*)$.

  \end{enumerate}

  $\adv$ wins the experiment if $\SPSverify(\pk, \vec{\msg}^*, \SPSsig^*) = 1$,
$|\Cor| \leq t$, and $|S(\vec{\msg}^*)| < t + 1 - |\Cor|$, where $S(\vec{\msg})$
records which honest signers have issued a partial signature on $\vec{\msg}$.

\end{experimentbox}

\paragraph{Security goals.} A threshold structure preserving signature scheme
must provide the following.

\begin{description}

  \item[Unforgeability.] No adversary corrupting at most $t$ signers can produce
a valid signature on a message vector for which it did not obtain partial
signatures from enough honest signers, so the scheme is no easier to forge
against than the structure preserving signature scheme underlying it. This is
\Cref{thm:tsps-construction-security}, proved in \Cref{sec:tsps-proof-1}.

  \item[Key secrecy.] No set of at most $t$ signers learns anything about the
signing key, at any point in the protocol, including key generation. This is
\Cref{lem:tsps-key-secrecy}, proved in \Cref{sec:tsps-proof-1}.

  \item[Transparency.] A public key produced by $\TSPSgen$ is distributed
exactly as one produced by $\SPSgen$, and a signature produced by $\TSPScombine$
verifies under $\SPSverify$ exactly as an ordinary signature would. Verification
is unchanged by thresholdisation. This is \Cref{lem:tsps-transparency}, proved
in \Cref{sec:tsps-proof-1}.

\end{description}

\subsection{Secure multiparty computation}\label{sec:tsps-mpc}

Let $\alpha \in \Zp$ be a secret value shared between the $n$ signers under a
linear $(t+1, n)$-threshold secret sharing scheme, so that any $t+1 $ of them
reconstruct it by a predefined linear combination and any $t$ of them learn
nothing about it. We write $\shr{\alpha}$ for such a sharing: the vector
$(\alpha_1, \ldots, \alpha_n)$ in which $\alpha_i$ is the share held by the
$i^{\mathrm{th}}$ signer. For $\vec{\alpha} = (\alpha_1, \ldots, \alpha_\ell)
\in \Zp^\ell$ we write $\shr{\vec{\alpha}}$ for $(\shr{\alpha_1}, \ldots,
\shr{\alpha_\ell})$, and we use the same notation $\shr{a}$ for a sharing of a
group element $a \in \group{}$, where the shares are themselves group elements
whose reconstruction is the corresponding product. We assume this notation
incorporates the authentication components that malicious security requires,
which we otherwise abstract away.

The reconstruction threshold and the corruption bound are related but distinct.
A dishonest majority realisation such as the one benchmarked in
\Cref{sec:tsps-evaluation} shares secrets additively between all $n$ parties,
which is the case $t = n-1$. every signer must then take part in reconstruction,
and the scheme is $n$-out-of-$n$ rather than $(t+1)$-out-of-$n$ for smaller $t$.
An honest majority realisation built on Shamir sharing supports any $t < n/2$
with genuine $(t+1)$-out-of-$n$ reconstruction. We state our constructions and
proofs for general $(t+1, n)$ and note the restriction where it bites.

Rather than specifying how the signers jointly compute each step of signing, we
describe our schemes as sequences of calls to a secure multiparty computation
protocol $\algo{MPC}$, and assume a realisation of it satisfying the properties
below. $\algo{MPC}$ is a resource our protocols consume, not a definitional
target they are proved to realise. We do not work in a hybrid model, and the
security statements of \Cref{sec:tsps-security} are made relative to these
properties rather than to any ideal functionality.

$\algo{MPC}$ is run jointly by the $n$ signers. Each interface below is invoked
by all of them together, and the stated output is the share each signer obtains.
Two of the interfaces operate on group elements, and each source group needs its
own copy of them, since a sharing over $\group{1}$ and a sharing over
$\group{2}$ are different objects and no interface mixes the two. We write
$\group{}$ when a statement holds for either group, and we subscript an
interface with $1$ or $2$ when the choice matters, so that $\algo{ExpGrp}_1$ is
the interface instantiated at $\group{1}$ and $\algo{ExpGrp}_2$ at $\group{2}$.
The interfaces over $\Zp$ are standard and are provided directly by existing MPC
frameworks. Those over $\group{}$ are not, and supplying them is what allows a
structure preserving signature to be thresholdised at all.

\begin{primitivebox}{Secure multiparty computation ($\algo{MPC}$)}{mpc}

  \begin{description}

    \item[$\shr{\alpha} \sample \algo{RandShare}()$ :] Outputs a $(t+1, n)$
sharing of a value $\alpha \sample \Zp$ that no signer learns.

    \item[$\shr{\gamma} \sample \algo{Mul}(\shr{\alpha}, \shr{\beta})$ :] On
input sharings of $\alpha$ and $\beta$, outputs a sharing of $\gamma =
\alpha\beta$.

    \item[$\shr{\gamma} \sample \algo{Inv}(\shr{\alpha})$ :] On input a sharing
of $\alpha \neq 0$, outputs a sharing of $\gamma = \alpha^{-1}$.

    \item[$\shr{c} \sample \algo{ExpGrp}(\shr{\alpha}, \shr{b})$ :] On input a
sharing of $\alpha \in \Zp$ and a sharing of $b \in \group{}$, outputs a sharing
of $c = b^\alpha$.

    \item[$\alpha \leftarrow \algo{Open}(\shr{\alpha})$ :] On input a sharing of
$\alpha$, reconstructs $\alpha$ from $t+1$ shares and distributes it to all
signers, or to a designated party. Applies equally to a sharing $\shr{a}$ of a
group element.

  \end{description}

\end{primitivebox}

\paragraph{Sharing group elements.} In practice $\group{}$ is an elliptic curve
group, whose elements are coordinate pairs $(x_0, y_0)$ satisfying the curve
equation. A group element is not shared by independently sharing its
coordinates, as $(\shr{x_0}, \shr{y_0})$, since computing a group operation on
such shares would require the signers to communicate. It is instead shared as a
single unit, $\shr{a} = \shr{(x_0, y_0)}$, with reconstruction the product of
the shares under the group law, so that the operations below remain local.  A
share of a group element may be produced without interaction from a share of its
discrete logarithm: given $\shr{\alpha}$, each signer computes $a_i =
\gen^{\alpha_i}$, and the resulting $\shr{a}$ shares $a = \gen^{\alpha}$. This
is how the group elements of a public key, and those parts of a secret key that
are group elements, are obtained in \Cref{sec:tsps-instantiation}.

\paragraph{Mixed-domain multiplication.} $\algo{ExpGrp}$ is the only interface
that combines the two domains, multiplying a secret-shared exponent by a secret
shared group element, and it is what makes $\algo{MPC}$ over $\group{}$ more
than a relabelling of $\algo{MPC}$ over $\Zp$. It is realised in the Beaver
style, consuming a mixed triple $(\shr{\eta}, \shr{h}, \shr{h^\eta})$ with $\eta
\sample \Zp$ and $h \sample \group{}$, generated during pre-processing.  The
signers open $\epsilon = \alpha - \eta$ over $\Zp$ and $\delta = b \cdot h^{-1}$
over $\group{}$, and each then computes its share of
%
\begin{equation*}
%
  c = b^{\alpha} = \left(h^{\eta}\right) \cdot \left(h^{\epsilon}\right) \cdot
\left(\delta^{\eta}\right) \cdot \delta^{\epsilon}
%
\end{equation*}
%
locally, from $\shr{h^\eta}$, $\shr{h}$, $\shr{\eta}$ and the two opened values,
with the public term $\delta^{\epsilon}$ folded in by a designated signer. The
masks $\epsilon$ and $\delta$ are uniform and independent of $\alpha$ and $b$,
so the two openings reveal nothing. Under malicious security the triple carries
authentication components in both domains, and it is supporting these over an
elliptic curve group that required the extension to MP-SPDZ described in
\Cref{sec:tsps-evaluation}.

\paragraph{Relation to existing work.} Computing on secret-shared group elements
is not new, and we set out here what we take from existing work and what we add.
That a sharing over $\Zp$ can be moved into a group without interaction, by each
signer raising a public generator to its own share, is folklore, and goes back
at least as far as verifiable secret sharing~\cite{FOCS:Feldman87}. Smart and
Talibi Alaoui give a full-threshold, actively secure protocol for computation
over an elliptic curve group. It takes Beaver triples produced over $\Zp$ in the
pre-processing model and maps them locally into triples over the curve, which is
the same route to $\algo{ExpGrp}$ that we take above~\cite{IMA:SmaTal19}.
Aranha, Dalskov, Escudero, and Orlandi give the most general treatment we are
aware of, under the name of linear secret sharing isomorphisms. These are linear
maps between sharings defined over vector spaces over a common field, and moving
from $\Zp$ into a source group is one instance of them. They extend the idea to
bilinear maps and hence to pairings~\cite{LC:ADEO21}. They apply it to
realise the Pointcheval-Sanders signature scheme in MPC, which is itself a
structure preserving scheme. Prior work has also shared group elements between
two parties and computed on those shares~\cite{NDSS:GPVRNC23}.

We therefore claim no novelty for the sharing $\shr{a}$ or for the idea of
computing on it, and we have stated the interfaces above in the form our
constructions consume rather than as a contribution in themselves. What this
chapter adds is the observation of \Cref{sec:tsps-construction}. The signing
algorithm of a structure preserving signature scheme decomposes so that its
message-dependent part uses only the local operations of the linearity
preservation property below. That is what yields an online phase requiring no
interaction at all. That, together with the two concrete schemes of
\Cref{sec:tsps-instantiation} and the extension to MP-SPDZ that makes them
runnable, is the contribution.

Two further lines of work are adjacent but concern different problems. Secret
sharing over abelian groups that are not vector spaces over a field is genuinely
harder, since the reconstruction coefficients can no longer be taken from a
field, and has its own literature~\cite{C:CraFeh02,C:DesFra91}.  Computation
over black-box, possibly non-abelian groups is harder still, as there is no
scalar action to exploit and even the $n$-fold product is
nontrivial~\cite{C:DPSW07,SCN:DesPieSte12}.

\paragraph{Inversion and non-zero sampling.} $\algo{Inv}$ is not a primitive
operation but is built from $\algo{Mul}$ and $\algo{Open}$: the signers sample
$\shr{\beta} \sample \algo{RandShare}()$, open $\gamma = \alpha\beta$, and take
$\shr{\alpha^{-1}} = \gamma^{-1}\shr{\beta}$, which is local. The opened
$\gamma$ is uniform and independent of $\alpha$, so nothing is revealed, and the
call costs one multiplication and one opening. Where a scheme calls for a value
sampled from $\Zp^*$ rather than $\Zp$, the signers cannot test a sharing for
zero without opening it. We instead sample from $\Zp$ and abort if the opening
performed inside $\algo{Inv}$ returns zero, which happens with probability $1/p$
and is therefore negligible. Where no inversion is performed, the difference
between the two distributions is statistically negligible and we elide it.

We require $\algo{MPC}$ to satisfy the following.

\paragraph{Linearity preservation.} Given sharings $\shr{\alpha}$ and
$\shr{\beta}$ and public values $\rho, \sigma, \tau \in \Zp$, the signers can
compute a sharing of $\gamma = \rho\alpha + \sigma\beta + \tau$ without
communication by each applying the combination to its own share:
%
\begin{equation*}
%
  \shr{\gamma} = \rho\shr{\alpha} + \sigma\shr{\beta} + \tau.
%
\end{equation*}
%
The same holds over $\group{}$: given sharings $\shr{a}$ and $\shr{b}$ and a
public group element $d$, the signers obtain $\shr{c} = \shr{a}^{\rho}
\shr{b}^{\sigma} d$ without communication. It extends to raising a public group
element to a shared exponent, $\shr{c} = d^{\shr{\alpha}}$, and to raising a
shared group element to a public exponent, $\shr{c} = \shr{a}^{\rho}$. None of
these require interaction, which is why $\algo{MPC}$ exposes no interface for
them. Since we assume malicious security, linearity must also be preserved for
the authentication components. It is this property, applied to the message
vector, that makes the online phase of our schemes non-interactive.

\paragraph{Privacy.} There is a PPT simulator $\algo{Sim}$ such that, for any
set $\Cor$ of at most $t$ corrupted signers and any sequence of calls to
$\algo{MPC}$, the joint view of $\Cor$ is simulatable from the corrupted
signers' inputs and outputs alone, together with any values disclosed by
$\algo{Open}$. In particular, the shares held by $\Cor$ are distributed
independently of every shared value the calls compute.

\paragraph{Security with abort.} Corrupt signers may supply shares of their own
choosing on input to any interface, and may cause it to return $\bot$ to every
honest signer, but cannot cause it to return an incorrect output.

\paragraph{Corruption model.} We require the guarantees above against an
adversary that corrupts signers adaptively, up to the bound $|\Cor| \leq t$,
since the unforgeability notion above is itself adaptive. Neither realisation
benchmarked in \Cref{sec:tsps-evaluation} is proved adaptively secure, so the
concrete schemes of \Cref{sec:tsps-instantiation} inherit the corresponding
static notion. \Cref{sec:tsps-proof-1} shows what adaptivity costs when it is
obtained by guessing instead.

\paragraph{Realisation.} Realising the interfaces of $\algo{MPC}$ over
$\group{}$ required us to extend MP-SPDZ with support for computation on
secret-shared group elements~\cite{CCS:Keller20}. How $\algo{MPC}$ is realised
is otherwise orthogonal to the constructions that follow.
\Cref{sec:tsps-evaluation} benchmarks two realisations, one assuming an honest
majority and one a dishonest majority.
